% TOP_PROBLEM: {"problemId": "problem.tuttes-3-flow-conjecture", "canonicalTitle": "Tutte's 3-Flow Conjecture", "releaseRank": 350, "primaryDomain": "combinatorics_discrete_geometry", "primaryDomainLabel": "Combinatorics & discrete geometry", "displayStatus": "open_with_solved_subcases", "releaseStatus": "open_with_solved_subcases"}
% !TeX program = lualatex
% Definition and sources only; this file is not an LLM solution attempt.
\documentclass[11pt]{article}
\usepackage[margin=25mm]{geometry}
\usepackage{fontspec}
\setmainfont{Latin Modern Roman}
\usepackage{amsmath,amssymb,mathtools}
\usepackage{unicode-math}
\setmathfont{Latin Modern Math}
\usepackage{xurl,hyperref}
\hypersetup{colorlinks=true,urlcolor=blue}
\setcounter{secnumdepth}{0}
\setlength{\parindent}{0pt}
\setlength{\parskip}{5pt}
\setlength{\emergencystretch}{3em}
\begin{document}
\section*{\texorpdfstring{Tutte's 3-Flow Conjecture}{Tutte's 3-Flow Conjecture}}\label{350}
\subsection{Definitions and mathematical statement}
A finite loopless multigraph is $4$-edge-connected if deleting any set of at most three edges leaves it connected. A nowhere-zero integer $3$-flow consists of an orientation and edge values $f(e)\in\{-2,-1,1,2\}$ obeying conservation at every vertex:
\[
 \sum_{e\text{ enters }v} f(e)=\sum_{e\text{ leaves }v}f(e).
\]
Tutte's conjecture says every such graph has this flow. Parallel edges are allowed, and the one-vertex loopless case has the empty flow. Negative values can equivalently be made positive by reversing their edges. The requirement is ordinary zero-boundary conservation; it does not ask that every prescribed admissible boundary modulo three can be realized.
\par\smallskip\textit{Formulation and concept references:} \href{https://arxiv.org/pdf/1406.1554}{[S1]}.

\subsection{Short English statement}
Does every four-edge-connected loopless multigraph admit a nonzero integer flow using only magnitudes one and two?
\subsection{Sources}
\begin{itemize}
\item[S1]Chen and Ning, A note on nowhere-zero 3-flow and Z3-connectivity, Definition1.1 and Conjecture1.2, PDF printed p2.
\url{https://arxiv.org/pdf/1406.1554}.
\textit{Formulation source.}
\end{itemize}
\end{document}
